\chapter{Sistemas de ecuaciones lineales}
\label{cap:sistemas}

Un \textbf{sistema lineal} de $m$ ecuaciones con $n$ incógnitas y
coeficientes en $\K$ se escribe
\[
 AX=B,\qquad A\in\M_{m\times n}(\K),\quad
 X=\begin{pmatrix}x_1\\\vdots\\x_n\end{pmatrix},\quad
 B=\begin{pmatrix}b_1\\\vdots\\b_m\end{pmatrix}.
\]
$A$ es la matriz de coeficientes y $[\,A\mid B\,]$, la
\textbf{matriz ampliada}. Una solución es un vector de valores que satisface
todas las ecuaciones. El sistema es \textbf{compatible} si tiene soluciones,
e \textbf{incompatible} si no las tiene. Es \textbf{homogéneo} si $B=0$;
en ese caso siempre admite la solución nula.

\section{Equivalencia y resolución por Gauss}
\label{sec:resolucion}

Dos sistemas con las mismas incógnitas son \textbf{equivalentes} si tienen
el mismo conjunto de soluciones. Intercambiar ecuaciones, sumar a una un
múltiplo de otra o multiplicarla por un escalar no nulo conserva las
soluciones. Estas operaciones se realizan por filas sobre la matriz
\emph{ampliada}, incluido el término independiente.
Una ecuación $0=0$ puede eliminarse; una ecuación $0=b$ con $b\ne0$
hace incompatible el sistema.

Al escalonar $[\,A\mid B\,]$ se distinguen dos situaciones:
\begin{enumerate}
 \item Hay una fila con todos los coeficientes nulos y término independiente
       no nulo: no existen soluciones.
 \item No hay contradicciones. Cada columna con pivote en la parte de
       coeficientes determina una \textbf{incógnita principal}; las demás
       son \textbf{libres o secundarias}. Se asignan parámetros arbitrarios
       a las libres y se despejan las principales desde la última ecuación.
\end{enumerate}
Con Gauss--Jordan, los despejes se leen directamente de la forma reducida.
Si hay $r$ pivotes en $A$, una solución compatible depende de $n-r$
parámetros libres. Al operar con parámetros en los coeficientes, deben
separarse los casos que anulan un pivote antes de dividir.

La forma reducida de la matriz ampliada es única. Para sistemas compatibles
con las mismas incógnitas, al suprimir las filas nulas se obtiene el mismo
sistema reducido si y solo si son equivalentes. En cambio, todos los
sistemas incompatibles tienen el mismo conjunto vacío de soluciones,
aunque sus formas reducidas puedan diferir.

\section{Discusión por rangos}
\label{sec:rouche}

El \textbf{teorema de Rouché--Frobenius} caracteriza la compatibilidad:
\[
 AX=B\text{ compatible}\quad\Longleftrightarrow\quad
 \rg(A)=\rg([\,A\mid B\,]).
\]
Es la misma condición que exigir que $B$ sea combinación lineal de
las columnas de $A$. Si $r=\rg(A)$,
\begin{center}
\begin{tabular}{@{}lll@{}}
\toprule
Condición & Tipo de sistema & Soluciones \\
\midrule
$r<\rg([\,A\mid B\,])$ & Incompatible & Ninguna. \\
$r=\rg([\,A\mid B\,])=n$ & Compatible determinado & Una. \\
$r=\rg([\,A\mid B\,])<n$ & Compatible indeterminado & Infinitas. \\
\bottomrule
\end{tabular}
\end{center}
Sobre $\R$ o $\C$, dos soluciones distintas ya producen infinitas:
si $X_1$ y $X_2$ resuelven el sistema, también lo hace
$\lambda X_1+(1-\lambda)X_2$ para todo $\lambda\in\K$.

Para el sistema homogéneo, el rango de la ampliada coincide con el de $A$.
Tiene solo la solución nula si $\rg(A)=n$, y soluciones no nulas si
$\rg(A)<n$. Su conjunto de soluciones es un subespacio de dimensión
$n-\rg(A)$, como se explica en la sección~\ref{sec:ecuaciones-subespacios}.

\section{Sistemas cuadrados y regla de Cramer}
\label{sec:cramer}

Si $A\in\M_n(\K)$, el sistema es compatible determinado si y solo si
$A$ es invertible. Su solución es $X=A^{-1}B$.
La \textbf{regla de Cramer}, válida cuando $\det(A)\ne0$, da
\[
 x_i=\frac{\det(A_i)}{\det(A)},\qquad i=1,\ldots,n,
\]
donde $A_i$ se obtiene sustituyendo la columna $i$ de $A$ por $B$.
La fórmula se deduce del producto de $B$ por la inversa expresada
mediante cofactores (sección~\ref{sec:determinantes}).

En un sistema compatible de rango $r>0$, puede elegirse un menor no nulo
de orden $r$. Sus filas seleccionan ecuaciones independientes y sus columnas,
$r$ incógnitas principales. Se pasan las otras $n-r$ incógnitas al segundo
miembro como parámetros. El sistema cuadrado restante es regular y permite
aplicar Cramer. Primero hay que verificar la compatibilidad del sistema
completo; el subsistema elegido por sí solo no la garantiza.
Si $r=0$, la compatibilidad exige $B=0$ y todas las incógnitas son libres;
no es necesario aplicar Cramer.

\section{Factorización LU}
\label{sec:lu}

Si una matriz regular admite una factorización $A=LU$, con $L$ triangular
inferior y $U$ triangular superior, ambas regulares, resolver $AX=B$
equivale a resolver sucesivamente
\[
 LY=B\quad\text{y}\quad UX=Y.
\]
El primer sistema se resuelve de arriba abajo y el segundo, de abajo arriba.
La factorización se reutiliza cuando cambia $B$ y se mantiene $A$.

\begin{samepage}
\subsection*{Existencia y construcción}

Para $A\in\M_n(\K)$ regular, son equivalentes:
\begin{itemize}
 \item Existe $A=LU$ con factores triangulares regulares.
 \item Gauss puede triangular $A$ sin intercambiar filas, en el orden habitual.
 \item $\Delta_k(A)\ne0$ para $k=1,\ldots,n$.
\end{itemize}
\end{samepage}
Aquí $\Delta_k$ se refiere al bloque de las primeras $k$ filas y columnas.
La regularidad de $A$ solo asegura $\Delta_n(A)\ne0$; los menores anteriores
también son necesarios para evitar intercambios.

Si $E_s\cdots E_1A=U$ mediante eliminaciones sin intercambios, entonces
$L=(E_s\cdots E_1)^{-1}$ es triangular inferior. Puede registrarse el
producto de operaciones triangularizando
\[
 [\,A\mid I_n\,]\longrightarrow[\,U\mid L^{-1}\,].
\]
Si se usan únicamente sumas de múltiplos para eliminar, $L$ tiene unos
en la diagonal y las entradas inferiores registran los multiplicadores
de eliminación con signo positivo.

Si se necesitan intercambios, una matriz de permutación $P$ permite
factorizar $PA=LU$. En ese caso se resuelven $LY=PB$ y luego $UX=Y$;
la permutación debe aplicarse también al término independiente.

\begin{ejemplo}
La factorización
\[
 \begin{pmatrix}2&1\\4&3\end{pmatrix}
 =\begin{pmatrix}1&0\\2&1\end{pmatrix}
  \begin{pmatrix}2&1\\0&1\end{pmatrix}
\]
se conserva para cualquier segundo miembro $B=(b_1,b_2)^t$.
Primero $y_1=b_1$, $y_2=b_2-2b_1$; después
$x_2=y_2$, $x_1=(y_1-x_2)/2$. Solo se repiten los dos despejes
triangulares cuando cambian los datos de $B$.
\end{ejemplo}
